\subsection{从端点一致性到轨迹映射与平均速度}
前面的自一致性条件本身还不够：如果网络总输出同一个常数，成对输出也会相等。
因此需要端点边界，例如 $f(x,0)=x$；
参数化中的 $c_{\rm skip}(0)=1,c_{\rm out}(0)=0$ 就用于落实这件事。
实际训练往往在一个小的正噪声 $\varepsilon$ 处设置边界，避免直接求零噪声分数。
此外，时间对的选取、离散化与距离函数都影响训练，
iCT 改进了独立一致性训练的这些设计，并区分了它与教师蒸馏时目标网络更新的作用。
不能把“两个输出相等”单独当作完整训练方案。\cite{cm,ict}

LCM 仍主要学习到固定低噪声端点的映射。
如果将目标时刻也作为输入，写成 $G(x,t,s)$，便可学习从 $t$ 到任意较低时刻 $s$ 的轨迹映射。
理想流映射满足
\[
G(x,t,t)=x,\qquad
G(G(x,t,r),r,s)=G(x,t,s),\quad s\leq r\leq t.
\]
第二式说的是，沿同一速度场先走一段再走一段，与直接走到目标时刻相同。
CTM 沿这个方向学习概率流轨迹，并配合相应训练约束。
这是把端点一致性扩展到两时刻之间的映射，并不是角色身份的一致性。\cite{ctm}

MeanFlow 则从平均速度继续推导。为了与前文扩散记号一致，
这里仍令 $t$ 向噪声端增加，生成时从 $t$ 走向 $s<t$。
对速度场 $v$ 的真实轨迹，定义
\[
u(x_t,s,t)=\frac1{t-s}\int_s^t v(x_r,r)\,\mathrm dr,
\qquad x_s=x_t-(t-s)u(x_t,s,t).
\]
这不是仅在起点求值的瞬时速度，而是整段轨迹的平均。
固定 $s$，沿轨迹对两边 $(t-s)u=\int_s^t v\,\mathrm dr$ 求导：
\[
u+(t-s)(\partial_tu+D_xu\,v)=v,
\]
从而得到
\[
u=v-(t-s)(\partial_tu+D_xu\,v).
\]
右边的方向导数可以用 Jacobian--vector product 计算，无需显式保存完整 Jacobian。
训练还需规定条件速度标签、停止梯度和权重等细节；
这个恒等式说明的是两时刻平均速度与瞬时速度的关系，
不是说把任意已训练的瞬时速度模型做一次大 Euler 步就能精确生成。\cite{meanflow,theorybook}
